% native creation stories: standards

\subsection{Hinduism}

There is no single account. The Rigveda gives several, and the tradition lets them stand together. One common shape runs through them.

Before anything, there was neither \textbf{asat} nor \textbf{sat}, neither non-being nor being. There was no sky, no air, no death, no deathlessness, no day, no night. There was only \textbf{Tad Ekam}, the One, which breathed without breath by its own impulse, depending on nothing outside itself. Beyond it there was nothing. Darkness lay hidden within darkness. Everything was \textbf{salila}, an undivided water with no shore.

Then something stirred. The Nasadiya Sukta names two forces. First \textbf{tapas}, the heat that drives one thing to separate from another. Then \textbf{kama}, desire, the first motion of mind reaching toward something other than itself. This kama was the first seed. With it, the undivided began to incline toward division. The sages, searching their own hearts, found there the bond that ties \textbf{sat} to \textbf{asat}.

Once the stirring began, the One started to become countable. The whole became the world.

The Purusha Sukta describes this as a \textbf{yajna}, a sacrifice. There was a vast cosmic person, \textbf{Purusha}, with a thousand heads and eyes and feet, who covered the whole earth and reached past it, who was all that had been and all still to come. The gods offered Purusha and divided him, and from the pieces the world was assembled. From his mind came the moon. From his eye came the sun. From his breath came Vayu the wind. From his head came the sky and from his feet the earth. Even the four \textbf{varnas} of human society came out of his body. Creation in this account is the whole broken into parts.

The Hiranyagarbha Sukta describes a golden womb. \textbf{Hiranyagarbha}, the golden embryo, arose first of all, floating on the formless water, holding within it the seed of everything that would be born. It became the one lord of all that came to be. It held up the earth and the sky and gave out life and breath. Every verse returns to the question \textbf{Kasmai devaya havisha vidhema}, to which god shall we offer. The hymn is a search for the one behind the many.

In the later Puranic account the egg, the \textbf{Brahmanda}, cracks open into the heavens and the earth, and from it \textbf{Brahma} is born, who orders the worlds, the beings, and the first \textbf{Prajapatis}. Brahma does not build from nothing. He re-forms what dissolved at the close of the previous round.

Creation is not a single event. It is one phase of a continuing rhythm. The worlds are breathed out, held for an immense span, then breathed back in. Three movements repeat under the \textbf{Trimurti}. A bringing forth, \textbf{srishti}, by Brahma at the dawn of each cosmic day. A long holding, \textbf{sthiti}, by Vishnu that keeps the worlds in order. And a dissolving, \textbf{samhara}, by Shiva at the close of each cosmic night, when everything returns to the unmanifest. At the end of a \textbf{kalpa} comes \textbf{pralaya}, and at the end of Brahma's life comes \textbf{mahapralaya}, when all returns to \textbf{Brahman}. After the night, the worlds form again.

The oldest hymn ends without certainty. After describing the beginning, it asks who could really know where creation came from. The gods came after the world was already made, so even they were not present to see it. Perhaps the One who watches from the highest heaven knows how it began, or perhaps even he does not know. The account is explicit that the first moment remains hidden.

\subsection{Christianity}

In the beginning there is a source that is pure goodness and beyond all telling, the \textbf{superessential} Godhead, the Good above being. It does not create out of need. It creates out of overflow, the way the sun gives light without losing any of itself. The source speaks, and the speaking is itself a power. By that \textbf{Logos}, the Word through whom all things are made, the worlds are made. Goodness passes downward through ordered ranks of being in the \textbf{proodos}, the procession, each rank receiving the light and passing it to the rank below. The same drawing, \textbf{eros}, that sends the light out also draws everything back in the \textbf{epistrophe}, the return. All things go out from the Good and all things are drawn home to it. The pattern is a circle, out and then back.

There is a second telling, the Gnostic account. In it the world we live in is not the work of the true source at all. At the summit is the unknowable \textbf{Monad}, also called \textbf{Bythos}, the Depth, beyond name and being. From it unfolds the \textbf{Pleroma}, the divine fullness, emanating in male-female pairs of \textbf{aeons}. The youngest aeon, \textbf{Sophia}, oversteps. She desires to grasp the source directly, outside the law of pairs. Her disordered act produces a flaw, a formless offspring cast outside the fullness. That offspring is the \textbf{Demiurge}, called \textbf{Yaldabaoth}, an ignorant maker born from the flaw above, who does not know the source over him. He mistakes himself for the highest and fashions the material world as a kind of prison, ruled by his \textbf{Archons} and by fate. Yet a divine spark of the true light is trapped inside the human being. Most people never notice it. A messenger comes from above to wake the spark, to remind it where it came from, and to show it the way back up past the Archons to the source. To attain \textbf{gnosis}, knowledge, is to be freed. The return of the scattered light to the Pleroma ends the account.

\subsection{Sufism}

In the beginning there is one \textbf{al-Haqq}, the Real. It is utterly one, alone, and beyond all likeness. Its \textbf{dhat}, the Essence, is the Hidden, never reached even by the prophets. In itself it is a hidden treasure, complete and unknown. A reported saying frames the whole. \textbf{Allah} says, I was a hidden treasure and I loved to be known, so I created the creation that I might be known. Out of that wish the world arises.

The one breathes its own \textbf{asma}, its names, outward, and each name takes form as a creature. Every thing that exists is the trace of a name, the outward showing of one face of the hidden one. The names split into two families. The \textbf{jamal} names disclose gentleness and nearness. The \textbf{jalal} names disclose majesty and rigor. So the cosmos is not a thing standing apart from the source. It is a mirror, a \textbf{mir'at}. In it the one sees its own names reflected back, and by seeing them it knows itself. This is \textbf{tajalli}, self-disclosure.

The breath that gives existence is \textbf{nafas al-Rahman}, the breath of the Merciful. To exist at all is to be touched by mercy. As speech is breath shaped by the mouth, the cosmos is the breath of God shaped by the names. This breathing is not a single past act. The world is re-breathed at every instant in \textbf{tajdid al-khalq}, made new moment by moment. Before things exist outwardly they exist as archetypes in the divine knowledge, the \textbf{a'yan thabita}, the fixed entities. They have no being of their own. They are the receptacles into which existence is poured. A thing brings its character from its own fixed entity.

Among all creatures the human being is the clearest part of the mirror, the \textbf{al-insan al-kamil}, the perfect man, the one being that can hold all the names at once. In the perfected human the reflection is complete, and the one beholds itself fully. The path of the soul is \textbf{fana}, to lose its separate self in the one, and then \textbf{baqa}, to abide in the one, so that the wave returns to the sea it never truly left.

\subsection{Sikhism}

For countless ages there was only formless dark. No earth, no sky, no day or night, no measure of time. There was only the one, alone and unmanifest, held in its own \textbf{hukam}, its will. The scripture opens with \textbf{Ik Onkar}, there is one. There is one supreme being, one \textbf{Karta Purakh}, the creator being, immanent in all and beyond all, timeless and self-existent, without fear and without hatred, never born and never dying. The devotional name for this one is \textbf{Waheguru}, the wondrous lord.

Then the one willed it, and with a single \textbf{Shabad}, a single word, brought forth the worlds. From that one utterance came the air, and from the air the water, and from the water the play of living things. The one set the worlds turning by its own \textbf{hukam} and poured itself into all of them. It is not a maker standing outside its work. It is present in every creature and every grain, the formless one within all forms.

The same will that brought the worlds out holds them in being and will draw them home again. The presence of the one in the heart is \textbf{Naam}, the divine name, and the practice of holding it in mind is \textbf{Naam simran}, remembrance. What blocks the soul is \textbf{haumai}, the I-am-ness, the sense of a separate self standing apart from the one. The whole of life is the unfolding of that one will, and the goal of the soul is \textbf{mukti}, liberation, the loss of haumai and the merging back into the one from which it came.

\subsection{Taoism}

In the beginning there is a nameless source. It cannot be spoken or pointed to. The opening of the \textbf{Tao Te Ching} states that the \textbf{Tao} that can be told is not the eternal Tao. It is silent and empty, formed in chaos before there was a sky or a ground, and it stands alone without changing. It is the mother of everything that will come. It shows two aspects that differ only in name, \textbf{Wu}, non-being, the nameless origin, and \textbf{You}, being, the named mother of the ten thousand things. At the threshold stand \textbf{Wuji}, the Limitless without distinction, and \textbf{Taiji}, the Supreme Ultimate, the first division.

Out of this source comes the \textbf{One}, a single undivided wholeness, the primordial \textbf{qi}, all potential held together in one breath. Chapter forty-two gives the formula. The Tao gives birth to the One. The One gives birth to the Two. The Two gives birth to the Three. The Three gives birth to the ten thousand things.

The One divides into the \textbf{Two}, a dark receiving side and a bright active side, \textbf{yin} and \textbf{yang}. Yin is the receptive, dark, cold, yielding, contracting aspect. Yang is the active, bright, warm, firm, expanding aspect. Neither is pure. Each holds a seed of the other. The two poles alone would fall back together, so a \textbf{Three} arises between them, the \textbf{chongqi}, a blending breath that holds them apart and joins them at once.

From this stable trio come the \textbf{wanwu}, the ten thousand things, the full crowd of the world, each carrying yin on its back and holding yang in its arms. The wise do not fight this flow. They practice \textbf{wu wei}, acting without forcing, moving like water that yields and still wears down stone. By acting this way they run the flow backward from the many to the one and return to the Tao, the nameless source.

\subsection{Hermeticism}

In the beginning there is a boundless Light, gentle and joyous, which is also \textbf{Nous}, Mind itself, God the Father, the unbegotten source. It is not a thing in the world. It is the ground of mind and being, the source from which everything flows. From it a part turns downward and away, condensing into \textbf{physis}, a dark, churning, watery Nature. The pure and the murky are now two, the above set against the below.

Then an ordering Word, the \textbf{Logos}, the holy Word and Son of God, descends upon the chaos and speaks the realms apart. The fiery and light rise, the heavy and moist settle below. From the source comes a second Mind, the \textbf{Nous-Demiourgos}, the craftsman, a god of fire and breath, who fashions the seven \textbf{dioiketores}, the Governors of the planetary spheres, that bind and govern the lower world. Their joined administration is \textbf{Heimarmene}, Fate, the turning of the spheres.

After this the source brings forth the \textbf{Anthropos}, the luminous archetypal Human, in its own image, equal to itself, free above the spheres and kin to the source. This Anthropos looks down through the spheres, sees its own form reflected in the water and the earth, is drawn to it, and descends. Nature embraces it, and the two unite. So the human becomes twofold, mortal in body yet immortal in essence. The descent comes not from sin imposed from outside but from love of one's own reflection. Then the androgynous beings are split into male and female, and reproduction begins.

The way back is knowledge. The self climbs back up through the circling spheres, shedding at each Governor the vice it imposed, until it passes the seventh gate, enters the \textbf{Ogdoad}, the eighth nature, sings praise with the powers above, surrenders into them, and itself becomes a power, entering into God. This is the Hermetic deification. The signature maxim of the \textbf{Tabula Smaragdina}, the Emerald Tablet, holds throughout. As above, so below.
